\documentclass[a4paper,11pt]{article}
\usepackage[margin=0.85in]{geometry}
\usepackage[T1]{fontenc}
\usepackage[utf8]{inputenc}
\usepackage{lmodern,microtype}
\usepackage[hidelinks]{hyperref}
\usepackage{enumitem}
\pagestyle{empty}
\setlength{\parindent}{0pt}
\setlength{\parskip}{8pt}
\begin{document}
\begin{flushright}
\textbf{Senanayak Sesh Kumar Karri}\\
Research Fellow, Imperial College London\\
\href{mailto:s.karri@imperial.ac.uk}{s.karri@imperial.ac.uk}\\
\href{https://seshkumar.github.io}{seshkumar.github.io}\\
14 October 2026
\end{flushright}
\textbf{Application for the G-Research NeurIPS 2026 travel and conference grant}

Dear Travel Grant Committee,

I am applying for support to attend NeurIPS 2026 in Sydney, where I have two papers accepted for poster presentation. I am a machine-learning researcher and Research Fellow at Imperial College London, based at the Brevan Howard Centre for Financial Analysis, Imperial College Business School. My research centres on probabilistic modelling, learning on graphs, neural-network geometry, and optimisation; financial markets provide an application domain for these methods.

My accepted papers contribute to core machine-learning methodology:
\begin{itemize}[leftmargin=*,itemsep=5pt,topsep=2pt]
\item \textbf{Practical Non-Stationary Graph Gaussian Processes} (first author; submission 25512). This work develops Gaussian-process models for heterogeneous graphs, linking node-level variability to interaction strength through an interpretable parameter. It addresses how probabilistic models can represent variation across a network rather than impose a uniform structure.
\item \textbf{Geometry-Constrained Kolmogorov--Arnold Networks: Learning Edge Geometry via Banach Duality} (sole author; \href{https://arxiv.org/abs/2608.25807}{arXiv:2608.25807}). This work derives learnable edge functions from Banach-space duality, allowing the geometry of each edge to adapt to data. It investigates interpretable representations and learning from noisy or limited observations.
\end{itemize}

I completed my PhD in machine learning at INRIA / \'Ecole Normale Sup\'erieure, working on probabilistic graphical models and submodular optimisation. Subsequent research appointments at IST Austria and Imperial's statistical machine-learning group and Data Science Institute developed this methodological focus. My publications include work at NeurIPS, ICML, JMLR, TMLR, and \textit{Machine Learning}. My current financial applications concern non-stationary time series, structured covariance estimation, and reinforcement learning for portfolio decisions.

Grant support would help meet the cost of travelling from London to Sydney and participating in the conference. Attendance would allow me to present both papers, obtain feedback from researchers in probabilistic modelling and neural architectures, and develop collaborations around these research questions.

Thank you for considering my application. My curriculum vitae accompanies this letter.

Yours sincerely,\\[3pt]
Senanayak Sesh Kumar Karri
\end{document}
